\documentclass{beamer}

\usepackage[T2A]{fontenc}
\usepackage[utf8]{inputenc}
\usepackage[english,russian]{babel}
\input{settings.tex}


\title{Теория Ляпунова}
\subtitle{Теория Управления, лекция 12}
\author{Сергей Савин}
\centering
\date{\mydate}


\begin{document}
\maketitle


%\begin{frame}{Content}
%
%\begin{itemize}
%\item Lyapunov method: stability criteria
%\item Lyapunov method: examples
%\item Linear case
%\item Discrete case
%\item Lyapunov equations
%\item Read more
%\end{itemize}
%
%\end{frame}




\begin{frame}{Метод Ляпунова: критерии устойчивости}
	\begin{flushleft}
		
		Рассмотрим автономную динамическую систему $\dot{\bo{x}} = \bo{f}(\bo{x})$, где $ \bo{f}(0) = 0$. Тогда:
		
		\begin{block}{Асимптотическая устойчивость}
			Данная система является асимптотически устойчивой, если существует скалярная функция $V = V(\bo{x}) > 0$, производная которой по времени отрицательна $\dot V(\bo{x}) < 0$, за исключением $V(\bo{0}) = 0$, $\dot V(\bo{0}) = 0$.
		\end{block}
		
		\begin{block}{Устойчивость по Ляпунову}
			$\dot{\bo{x}} = \bo{f}(\bo{x})$ устойчива в смысле Ляпунова, если $\exists V(\bo{x}) > 0$, $\dot V(\bo{x}) \leq 0$.
		\end{block}
		
		\begin{definition}
			Функция $V(\bo{x}) > 0$, удовлетворяющая этим условиям, называется \emph{функцией Ляпунова}.
		\end{definition}
		
	\end{flushleft}
\end{frame}

\begin{frame}{Метод Ляпунова: Пример 1}
	\begin{flushleft}
		
		Рассмотрим динамическую систему $\dot{x} = -x$. 
		
		\bigskip
		
		Предложим \emph{кандидата на функцию Ляпунова} $V(x) = x^2 \geq 0$. Найдём её производную:
		
		\begin{equation}
			\dot V(x) = 
			\frac{d}{dt}(  x^2  ) = 
			2x \dot x =
			2x (-x) =
			-x^2 \leq 0
		\end{equation}
		
		Это удовлетворяет условиям асимптотической устойчивости.
		
	\end{flushleft}
\end{frame}

\begin{frame}{Метод Ляпунова: Пример 2}
	\begin{flushleft}
		
		Рассмотрим маятник $\ddot{q} = f(q, \dot{q}) = -\dot{q} - \sin(q)$. 
		
		\bigskip
		
		Предложим \emph{кандидата на функцию Ляпунова} $V(q, \dot{q}) = E(q, \dot{q}) = \frac{1}{2} \dot{q}^2 + 1 - \cos(q)\geq 0$, где $E(q, \dot{q})$ — полная энергия системы. Найдём её производную:
		
		
		\begin{align}
			\dot V(q, \dot{q}) = 
			\frac{d}{dt} (\frac{1}{2} \dot{q}^2 + 1 - \cos(q)) =
			\dot{q} \ddot{q} + \sin(q)\dot{q} = \\
			=
			\dot{q} (-\dot{q} - \sin(q)) + \sin(q)\dot{q}
			=
			-\dot{q}^2 \leq 0
		\end{align}
		
		Это удовлетворяет критериям Ляпунова, поэтому система устойчива. Асимптотическая устойчивость не доказана, поскольку $\dot V(q, \dot{q}) = 0$ для всех $(q, \dot q)$ таких, что $\dot q = 0$.
		
	\end{flushleft}
\end{frame}

\begin{frame}{Принцип инвариантности ЛаСалля, 1}
	\begin{flushleft}
		
		\begin{block}{Принцип инвариантности ЛаСалля}
			Автономная динамическая система $\dot{\bo{x}} = \bo{f}(\bo{x})$ является асимптотически устойчивой, если существует скалярная функция $V = V(\bo{x}) > 0$, производная которой по времени неположительна $\dot V(\bo{x}) \leq 0$, за исключением $V(\bo{0}) = 0$, где множество $\{\bo{x}: \   \dot V(\bo{x}) = 0  \}$ не содержит нетривиальных траекторий.
		\end{block}
		
		\bigskip
		
		Тривиальная траектория — это $\bo{x}(t) = 0$. В отличие от условия Ляпунова, принцип ЛаСалля позволяет доказывать асимптотическую устойчивость даже для систем с $\dot V(\bo{x}) = 0$.
		
	\end{flushleft}
\end{frame}

\begin{frame}{Принцип инвариантности ЛаСалля, 2}
	\begin{flushleft}
		
		Локальная версия принципа инвариантности ЛаСалля имеет следующий вид:
		
		\begin{block}{Локальный принцип инвариантности ЛаСалля}
			Автономная динамическая система $\dot{\bo{x}} = \bo{f}(\bo{x})$ является асимптотически устойчивой в окрестности $\mathcal D$ начала координат, если существует скалярная функция $V = V(\bo{x}) > 0$ (за исключением $V(\bo{0}) = 0$), производная которой по времени неположительна $\dot V(\bo{x}) \leq 0$, где множество $\mathcal M = \{\bo{x}: \   \dot V(\bo{x}) = 0  \} \cap \mathcal D$ не содержит нетривиальных траекторий.
		\end{block}
		
	\end{flushleft}
\end{frame}

\begin{frame}{Принцип ЛаСалля: Пример 2}
	\begin{flushleft}
		
		В нашем предыдущем примере $\dot V(q, \dot{q}) = 0$ для любого $q$, пока $\dot{q} = 0$. Но множество $\{(q, \dot{q}): \   \dot{q} = 0 \}$ не содержит траекторий системы $\ddot{q}  = -\dot{q} - \sin(q)$, кроме $q(t) = 0$, в области $-\frac{\pi}{2} < q < \frac{\pi}{2}$. Следовательно, принцип ЛаСалля доказывает локальную асимптотическую устойчивость.
		
	\end{flushleft}
\end{frame}

\begin{frame}{Принцип ЛаСалля: Пример 3}
	\begin{flushleft}
		
		Рассмотрим осциллятор $\ddot{q} = f(q, \dot{q}) = -\dot{q}$. 
		
		\bigskip
		
		Предложим \emph{кандидата на функцию Ляпунова} $V(q, \dot{q}) = T(q, \dot{q}) = \frac{1}{2} \dot{q}^2 \geq 0$, где $T(q, \dot{q})$ — кинетическая энергия системы. Найдём её производную:
		
		\begin{equation}
			\dot V(q, \dot{q}) = 
			\frac{\partial V}{\partial q}       \dot{q} +
			\frac{\partial V}{\partial \dot{q}} f(q, \dot{q}) = 
			\dot{q} (-\dot{q}) = -\dot{q}^2 \leq 0
		\end{equation}
		
		Это удовлетворяет критериям Ляпунова, поэтому система устойчива. Заметим, что $\dot V(q, \dot{q}) = 0$ для любого $q$, пока $\dot{q} = 0$. Но множество $\{(q, \dot{q}): \   \dot{q} = 0 \}$ содержит бесконечно много траекторий системы $\ddot{q} = -\dot{q}$, отличных от $q(t) = 0$, например $q(t) = 1$ или $q(t) = -2$. Следовательно, принцип ЛаСалля в этом случае не доказывает асимптотическую устойчивость.
		
	\end{flushleft}
\end{frame}

\begin{frame}{Линейный случай, 1}
	\begin{flushleft}
		
		Как вы видели, метод Ляпунова позволяет работать как с нелинейными, так и с линейными системами. Но для линейных систем существуют дополнительные свойства, которые можно использовать.
		
		\bigskip
		
		\begin{block}{Наблюдение 1}
			Для линейной системы $\dot{\bo{x}} = \bo{A}\bo{x}$ мы всегда можем выбрать кандидата на функцию Ляпунова в виде $V = \bo{x}^\top\bo{S}\bo{x} > 0$, где $\bo{S}$ — положительно определённая матрица.
		\end{block}
		
		\bigskip
		
		Следующий слайд покажет, к чему это приводит.
		
	\end{flushleft}
\end{frame}

\begin{frame}{Линейный случай, 2}
	\begin{flushleft}
		
		Для $\dot{\bo{x}} = \bo{A}\bo{x}$ и $V = \bo{x}^\top\bo{S}\bo{x} > 0$ найдём её производную:
		
		\begin{equation}
			\dot V(\bo{x}) = \dot{\bo{x}}^\top\bo{S}\bo{x} + 
			\bo{x}^\top\bo{S}\dot{\bo{x}} < 0
		\end{equation}
		
		\begin{equation}
			\dot V(\bo{x}) = (\bo{A}\bo{x})^\top\bo{S}\bo{x} + 
			\bo{x}^\top\bo{S}\bo{A}\bo{x} = 
			\bo{x}^\top(\bo{A}^\top\bo{S} + \bo{S}\bo{A})\bo{x} < 0
		\end{equation}
		
		Заметим, что $\dot V(x)$ должна быть отрицательной для всех $\bo{x}$, чтобы система была асимптотически устойчивой, что означает, что $\bo{A}^\top\bo{S} + \bo{S}\bo{A}$ должна быть отрицательно определённой. Другая форма этого требования — \emph{уравнение Ляпунова}:
		
		\begin{equation}
			\bo{A}^\top\bo{S} + \bo{S}\bo{A} = -\bo{Q}
		\end{equation}
		
		где $\bo{Q}$ — положительно определённая матрица.
		
	\end{flushleft}
\end{frame}

\begin{frame}{Дискретный случай, 1}
	\begin{flushleft}
		
		\begin{block}{Критерий асимптотической устойчивости, дискретный случай}
			Для системы $\bo{x}_{i+1} = \bo{f}(\bo{x}_i)$, если $V(\bo{x}_i) > 0$ и $V(\bo{x}_{i+1}) - V(\bo{x}_i) < 0$, то система асимптотически устойчива.
		\end{block}
		
		\bigskip 
		
		Как и прежде, для линейных систем мы будем выбирать \emph{положительно определённые квадратичные формы} в качестве кандидатов на функцию Ляпунова.
		
	\end{flushleft}
\end{frame}

\begin{frame}{Дискретный случай, 2}
	\begin{flushleft}
		
		Рассмотрим динамику $\bo{x}_{i+1} = \bo{A}\bo{x}_i$ и $V = \bo{x}_i^\top\bo{S}\bo{x}_i > 0$, найдём $V(\bo{x}_{i+1}) - V(\bo{x}_i)$:
		
		\begin{align}
			V(\bo{x}_{i+1}) - V(\bo{x}_i) &= \bo{x}_{i+1}^\top\bo{S}\bo{x}_{i+1}- 
			\bo{x}_i^\top\bo{S}\bo{x}_i =
			\\
			&= (\bo{A}\bo{x}_i)^\top\bo{S}\bo{A}\bo{x}_i - 
			\bo{x}_i^\top\bo{S}\bo{x}_i
			\\
			&= \bo{x}_i^\top(\bo{A}^\top\bo{S}\bo{A} - \bo{S})\bo{x}_i < 0
		\end{align}
		
		Заметим, что $V(\bo{x}_{i+1}) - V(\bo{x}_i)$ должна быть отрицательной для всех $\bo{x}_i$, чтобы система была асимптотически устойчивой, что означает, что $\bo{A}^\top\bo{S}\bo{A} - \bo{S}$ должна быть отрицательно определённой, что даёт нам \emph{дискретное уравнение Ляпунова}:
		
		\begin{equation}
			\bo{A}^\top\bo{S}\bo{A} - \bo{S} = -\bo{Q}
		\end{equation}
		
		где $\bo{Q}$ — положительно определённая матрица.
		
	\end{flushleft}
\end{frame}

\begin{frame}{Уравнения Ляпунова}
	\begin{flushleft}
		
		На практике вы можете легко использовать уравнения Ляпунова для проверки устойчивости. В Python и MATLAB есть встроенные функции для их решения:
		
		\begin{itemize}
			\item scipy: \texttt{linalg.solve\_continuous\_lyapunov(A, Q)}
			\item MATLAB: \texttt{lyap(A,Q)}
		\end{itemize}
		
	\end{flushleft}
\end{frame}






\begin{frame}{Литература}
\begin{flushleft}

\begin{itemize}
    \item \bref{https://folk.uib.no/nmagb/m2142002l3.pdf}{3.9 Liapunov’s direct method}
    \item \bref{https://arxiv.org/abs/1809.05289}{Universita degli studi di Padova Dipartimento di Ingegneria dell'Informazione, Nicoletta Bof, Ruggero Carli, Luca Schenato, Technical Report, Lyapunov Theory for Discrete Time Systems}
\end{itemize}

\end{flushleft}
\end{frame}




\myqrframe

\end{document}
